% !TEX TS-program = xelatex
% pt-PT, AO45, UTF-8
\lccode`\-=`\-
\documentclass[12pt,a4paper,twoside]{report}
% PREÂMBULO REORGANIZADO E COMENTADO — UTF-8 — pt-PT, AO45
% Uso: depois de documentclass, escrever input{auxiliares/preambulo.tex}.
% Motor previsto: XeLaTeX (xcolor e hyperref conservam a opção xetex).
%
% PRESERVAÇÃO E DEPENDÊNCIAS
% Todos os pacotes, opções, macros e valores do original foram conservados.
% amsmath passa a preceder unicode-math; hyperref precede o código que usa href.
% As fontes EB Garamond e NewCMMath-Book.otf têm de estar disponíveis.
% bibliografia.bib deve existir no projecto; a bibliografia requer Biber.
% logos/aeB05.pdf é usado nos cabeçalhos a partir da terceira página.
% O ambiente codigo usa number within=chapter: exige uma classe com chapter,
% como book ou report. Numa classe article, a opção tem de ser adaptada;
% conserva-se aqui a configuração original, sem mudar a numeração pretendida.
%
%=======================================================================
% 01. CORES BÁSICAS E FUNDO DE PÁGINA
%=======================================================================
% xcolor: Define cores, mistura-as e permite colorir texto, tabelas e o fundo da
% página. As opções dvipsnames e svgnames disponibilizam nomes de cores; table activa
% as cores nas tabelas.
\usepackage[xetex,dvipsnames,svgnames,table]{xcolor}

% Cores definidas por CMYK
% \definecolor{fondpaille}{cmyk}{0,0.02,0.2,0}
\definecolor{fondpaille}{cmyk}{0, 0.0118, 0.1176, 0.01} %cmyk(0%, 2.75%, 23.53%, 0%, 1%)
\definecolor{ferrugem}{cmyk}{0, 0.4457, 0.75, 0.6392}

% Cores definidas por hexadecimal
\definecolor{moleskin}{HTML}{FFF8DC}
\definecolor{laranja}{HTML}{F07241}
\definecolor{marron}{HTML}{800000}
\definecolor{lovely}{HTML}{C04848}
\definecolor{click}{HTML}{480048}
\definecolor{tinta-velha}{HTML}{4e4643}
\definecolor{Browned_Cherry}{HTML}{910831}
\definecolor{azul}{HTML}{000080}

% Cores da capa
\definecolor{BLACK}{HTML}{000000}
\definecolor{ELECTRICRED}{HTML}{DD0000}
\definecolor{TANGERINYELLOW}{HTML}{FFCE00}

% Cor consolidada para texto principal
\colorlet{cor-da-tinta}{ferrugem}
\newif\ifferrugem
\ferrugemfalse

% Cor de fundo da página
\pagecolor{fondpaille}

%=======================================================================
% 02. FONTES
%=======================================================================
% fontspec: Selecciona fontes OpenType e as suas características tipográficas com
% XeLaTeX ou LuaLaTeX.
\usepackage{fontspec} % XeLaTeX/LuaLaTeX fonts
% Mantêm-se todos os recursos OpenType originais: algarismos antigos,
% ligaduras e variantes. O sinal + activa um recurso; o sinal - desactiva-o.
\setmainfont{EBGaramond}[ % Tudo junto, não 'EB Garamond'
  %%  %%  %%  %%  %%  %%  %%  %%  %%
  Path=./auxiliares/ebgaramond/,
  Extension=.ttf,
  UprightFont=*-Regular,
  BoldFont=*-Bold,
  ItalicFont=*-Italic,
  BoldItalicFont=*-BoldItalic,
  %%  %%  %%  %%  %%  %%  %%  %%  %%
  Numbers=OldStyle,
  RawFeature={
    -aalt,
    -c2pc,
    -c2sc,
    -case,
    +dlig,
    -frac,
    -hist,
    +hlig,
    +kern,
    +liga,
    +mark,
    +mkmk,
    +onum,
    +pnum,
    +ss05,
    +cv06
  }
]
% amsmath é carregado antes de unicode-math para evitar alterações tardias.
% amsmath: Fornece alinhamentos, equações em várias linhas, intertext e outros
% recursos de composição matemática.
\usepackage{amsmath}
% unicode-math: Usa fontes matemáticas Unicode; math-style=ISO e bold-style=ISO
% escolhem as convenções ISO para os símbolos.
\usepackage[math-style=ISO, bold-style=ISO]{unicode-math}
% \setmathfont{OldStandard-Math.otf}
% \setmathfont{Garamond-Math.otf}[StylisticSet={5,6,7,9}]
\setmathfont{NewCMMath-Book.otf}

%=======================================================================
% 03. LÍNGUAS E HIFENIZAÇÃO
%=======================================================================
% polyglossia: Define as línguas, a hifenização e as designações automáticas do
% documento. A opção de língua não revê a ortografia: o texto deve ser escrito em
% pt-PT, AO45.
\usepackage{polyglossia}
\setdefaultlanguage{portuguese} %pt-PT, AO45
\PolyglossiaSetup{portuguese}{indentfirst=true}
\setotherlanguages{english,french}  % Apenas os que precisas

%=======================================================================
% 04. LAYOUT DE PÁGINA
%=======================================================================
% geometry: Define o formato e as margens da página; includehead e includefoot incluem
% cabeçalho e rodapé no cálculo da área.
\usepackage[a4paper,left=2.5cm,right=2.5cm,top=2.5cm,bottom=2.5cm,includehead,includefoot]{geometry}
% setspace: Disponibiliza comandos para alterar o espaçamento entre linhas.
\usepackage{setspace} % espaçamento de linhas
% lettrine: Cria capitulares que ocupam várias linhas no início de um parágrafo.
% calligra: Fornece uma fonte caligráfica, usada aqui nas capitulares.
\usepackage{lettrine,calligra} % letras capitulares
\renewcommand{\LettrineFontHook}{\calligra}
\LettrineTextFont{\itshape}
\setcounter{DefaultLines}{3}

% Prevenção de viúvas e órfãos
\clubpenalty=100
\linepenalty=100
\widowpenalty=100

% Outros ajustes tipográficos
% xurl: Melhora as oportunidades de quebra de linha dos endereços URL.
\usepackage{xurl}
% caption: Configura legendas; hang usa composição suspensa, small reduz o tamanho e
% bf selecciona negrito para a etiqueta.
\usepackage[hang,small,bf]{caption}
\setlength{\captionmargin}{1.5in}
\renewcommand{\fboxrule}{2pt}

%=======================================================================
% 05. TIPOGRAFIA AVANÇADA
%=======================================================================
% microtype: Melhora a composição do texto com ajustes tipográficos suportados pelo
% motor e pelas fontes.
\usepackage{microtype}      % ajustes de tipografia
% footmisc: Modifica as notas de rodapé; bottom coloca-as no fundo da área disponível
% da página.
\usepackage[bottom]{footmisc} % pés de página alinhados ao fundo
% \linespread{1.3}            % espaçamento de linha
% \linespread{1.2/1.25/1.3}

%=======================================================================
% 06. MATEMÁTICA E SÍMBOLOS
%=======================================================================
% dsfont: Fornece letras matemáticas com traço duplo através de mathds.
% mathrsfs: Fornece um alfabeto caligráfico matemático através de mathscr.
\usepackage{dsfont,mathrsfs}
% cancel: Permite riscar termos matemáticos; makeroom reserva espaço e thicklines
% engrossa os traços.
\usepackage[makeroom,thicklines]{cancel} % para riscar
% esvect: Desenha setas de vectores através de vv; conserva-se para os documentos que
% usam esse comando.
\usepackage{esvect}       % vetores

%=======================================================================
% 07. DATAS E FORMATAÇÃO NUMÉRICA
%=======================================================================
% datetime: Formata datas e horas; conservam-se as opções originais nodayofweek e level.
\usepackage[nodayofweek,level]{datetime}
% siunitx: Compõe números e unidades; output-decimal-marker selecciona a vírgula decimal.
\usepackage[output-decimal-marker={,}]{siunitx}

%=======================================================================
% 08. PACOTES DE CÁLCULO E EXERCÍCIOS
%=======================================================================
% xlop: Apresenta operações aritméticas e cálculos em disposição tradicional.
\usepackage{xlop}      % cálculo automático

%=======================================================================
% 09. GRÁFICOS E FIGURAS
%=======================================================================
% graphicx: Insere imagens e permite alterar a sua escala, rotação e recorte.
\usepackage{graphicx}
% svg: Insere imagens SVG; a conversão pode requerer Inkscape e autorização para
% executar comandos externos.
\usepackage{svg}
% tikz: Cria desenhos vectoriais, esquemas e elementos sobrepostos à página.
\usepackage{tikz}
% tkz-fct: Acrescenta recursos para desenhar funções; a execução de alguns recursos
% pode depender de ferramentas externas.
\usepackage{tkz-fct}
% tikzpagenodes: Disponibiliza nós TikZ correspondentes às áreas da página.
\usepackage{tikzpagenodes}
% circuitikz: Desenha circuitos eléctricos e os seus componentes.
\usepackage{circuitikz}
% pgfplots: Desenha gráficos de funções e dados sobre TikZ.
% pgfplotstable: Lê, transforma e apresenta dados em tabelas, em articulação com pgfplots.
\usepackage{pgfplots,pgfplotstable}
\pgfplotsset{compat=1.18}
\usetikzlibrary{
  calc,
  positioning,
  automata,
  arrows,
  shadows,
  shapes.multipart,
  matrix,
  backgrounds,
  decorations.pathreplacing,
  fit,
  plotmarks,
  tikzmark
}
\pgfdeclarelayer{background}
\pgfdeclarelayer{foreground}
\pgfsetlayers{background,main,foreground}

%=======================================================================
% 10. TABELAS AVANÇADAS
%=======================================================================
% array: Amplia a definição e a formatação das colunas das tabelas.
% longtable: Permite tabelas com continuação automática entre páginas e cabeçalhos
% repetidos.
% booktabs: Fornece regras horizontais adequadas à composição de tabelas, como top-,
% mid- e bottomrule.
% bigstrut: Acrescenta suportes invisíveis para aumentar a altura das linhas de tabelas.
% tabularx: Ajusta as colunas X à largura total indicada; o ambiente tabularx, por si
% só, não se divide entre páginas.
% seqsplit: Permite quebrar sequências longas de caracteres através de seqsplit.
\usepackage{array,longtable,booktabs,bigstrut,tabularx,seqsplit}

%=======================================================================
% 11. UTILITÁRIOS DIVERSOS
%=======================================================================
% comment: Permite excluir blocos inteiros através do ambiente comment.
\usepackage{comment}
% enumerate: Permite personalizar as etiquetas das listas enumerate. Não fornece a
% sintaxe label= e start= do pacote enumitem.
\usepackage{enumerate}
\newcounter{lastenum}
\newcommand{\storelast}{\setcounter{lastenum}{\value{enumi}}}
\newcommand{\uselast}{\setcounter{enumi}{\value{lastenum}}}
\newcommand\CC{C\nolinebreak[4]\hspace{-.05em}\raisebox{.4ex}{\relsize{-3}{\textbf{++}}}}
% rotating: Permite rodar objectos e criar figuras ou tabelas em orientação lateral.
\usepackage{rotating}
% lastpage: Cria uma referência à última página, útil para indicar o total de páginas.
\usepackage{lastpage}
% refcount: Permite obter valores numéricos a partir de referências.
\usepackage{refcount}
% standalone: Permite incorporar documentos autónomos e aproveitar os respectivos
% conteúdos.
\usepackage{standalone}
% relsize: Altera o tamanho da letra relativamente ao tamanho actual; é usado na macro CC.
\usepackage{relsize}
% ulem: Fornece sublinhados e outros realces; normalem preserva a forma habitual de emph.
\usepackage[normalem]{ulem}

%=======================================================================
% 12. SECÇÕES E SUMÁRIO
%=======================================================================
% titlesec: Configura títulos e espaçamentos e permite definir níveis adicionais de
% secção.
\usepackage{titlesec}
% tocloft: Configura a apresentação do sumário e das listas de figuras e tabelas.
\usepackage{tocloft}
% etoolbox: Fornece comandos de programação, testes e alterações de macros.
\usepackage{etoolbox}

% Novos níveis de secção
\titleclass{\subsubsubsection}{straight}[\subsubsection]
\titleclass{\subsubsubsubsection}{straight}[\subsubsubsection]
\newcounter{subsubsubsection}[subsubsection]
\newcounter{subsubsubsubsection}[subsubsubsection]
\renewcommand{\thesubsubsubsection}{\thesubsubsection.\arabic{subsubsubsection}}
\renewcommand{\thesubsubsubsubsection}{\thesubsubsubsection.\arabic{subsubsubsubsection}}

% Formato visual
\titleformat{\subsubsubsection}{\normalfont\normalsize\bfseries}{\thesubsubsubsection}{1em}{}
\titlespacing*{\subsubsubsection}{0pt}{3.25ex plus 1ex minus .2ex}{1.5ex plus .2ex}
\titleformat{\subsubsubsubsection}{\normalfont\normalsize\bfseries}{\thesubsubsubsubsection}{1em}{}
\titlespacing*{\subsubsubsubsection}{0pt}{3.25ex plus 1ex minus .2ex}{1.5ex plus .2ex}

% Sumário detalhado
\renewcommand{\cftsecfont}{\bfseries}
\renewcommand{\cftsubsecfont}{\normalfont}
\renewcommand{\cftsubsubsecfont}{\normalfont}
\makeatletter
\renewcommand{\l@section}{\@dottedtocline{1}{1.5em}{1.5em}}
\renewcommand{\l@subsection}{\@dottedtocline{2}{2.5em}{2em}}
\renewcommand{\l@subsubsection}{\@dottedtocline{3}{4em}{2.5em}}
\newcommand{\l@subsubsubsection}{\@dottedtocline{4}{6em}{3em}}
\newcommand{\l@subsubsubsubsection}{\@dottedtocline{5}{8em}{4em}}
\providecommand*{\toclevel@subsubsubsection}{4}
\providecommand*{\toclevel@subsubsubsubsection}{5}
\makeatother
\setcounter{secnumdepth}{5}
\setcounter{tocdepth}{5}

%=======================================================================
% 13. CÓDIGO-FONTE E PYTHON
%=======================================================================
% Carregar fvextra (e a sua dependência lineno) antes de csquotes.
% Esta ordem evita o aviso de fvextra e a alteração de @parboxrestore.
% fvextra: Amplia os ambientes verbatim de fancyvrb, incluindo recursos de quebra de
% linhas.
\usepackage{fvextra}
% fancyvrb: Apresenta texto e código literalmente, sem interpretar os seus caracteres
% como comandos LaTeX.
\usepackage{fancyvrb}

% pythontex: Integra código Python no documento e permite incorporar resultados. A
% execução requer o passo externo PythonTeX, além da compilação LaTeX.
% \usepackage{pythontex}

%=======================================================================
% 14. NOTAS DE RODAPÉ E CITAÇÕES
%=======================================================================
% fvextra já foi carregado na secção anterior, antes de csquotes.
% csquotes: Fornece citações e aspas adaptadas à língua, com integração em biblatex.
\usepackage{csquotes}
% appendix: Organiza os apêndices; toc e page permitem a sua entrada no sumário e uma
% página própria.
\usepackage[toc,page]{appendix}
\renewcommand\appendixpagename{Apêndices}
% biblatex: Gere as referências e citações bibliográficas. Este documento usa Biber,
% citações numéricas compactas e a ordem de citação.
\usepackage[
  backend=biber,
  style=numeric-comp,
  language=auto,
  natbib=false,
  block=space,
  isbn=false,
  url=true,
  doi=true,
  mcite=true,
  eprint=true,
  sorting=none
]{biblatex}
\addbibresource{bibliografia.bib}

% As notas normais usam letras; a marca é apresentada entre parênteses.
% symbolfootnote permite escolher uma marca simbólica para uma nota.
% Formato das notas de rodapé
\renewcommand{\thefootnote}{\alph{footnote}}
\makeatletter
\renewcommand\@makefnmark{\@textsuperscript{\normalfont(\@thefnmark)}}
\makeatother
\long\def\symbolfootnote[#1]#2{%
  \begingroup
  \def\thefootnote{\fnsymbol{footnote}}
  \footnote[#1]{#2}
  \endgroup
}

%=======================================================================
% 15. CÓDIGO E CAIXAS (QUEBRÁVEIS ENTRE PÁGINAS)
%=======================================================================
% listings: Apresenta código-fonte com realce de sintaxe, numeração de linhas e
% estilos; não executa o código.
\usepackage{listings}
% tcolorbox: Cria caixas com títulos, cores e molduras; most carrega um conjunto amplo
% de bibliotecas.
\usepackage[most]{tcolorbox} % carregar UMA SÓ VEZ no preâmbulo
\tcbuselibrary{listings,breakable}

%=======================================================================
% 16. CORES (CÓDIGO)
%=======================================================================
\definecolor{codebg}{gray}{0.96}
\definecolor{codeframe}{gray}{0.75}
\definecolor{codeblue}{RGB}{40,40,180}
\definecolor{codegreen}{RGB}{0,120,0}
\definecolor{codered}{RGB}{160,20,20}
\definecolor{darkgray}{gray}{0.4}

%=======================================================================
% 17. TRADUÇÃO LISTING → LISTAGEM
%=======================================================================
\renewcommand{\lstlistingname}{Listagem}
\renewcommand{\lstlistlistingname}{Lista de Listagens}

%=======================================================================
% 18. ESTILO LISTINGS (ESTÁVEL COM XELATEX)
%=======================================================================
% Definir o estilo não o aplica a todos os ambientes lstlisting.
% Usar [style=codigo] num lstlisting ou o ambiente codigo definido abaixo.
\lstdefinestyle{codigo}{
  language=Python,
  basicstyle=\ttfamily\small,
  keywordstyle=\color{codeblue}\bfseries,
  commentstyle=\color{codegreen}\itshape,
  stringstyle=\color{codered},
  numbers=left,
  numberstyle=\tiny\color{darkgray},
  numbersep=8pt,
  breaklines=true,
  showstringspaces=false,
  columns=fullflexible,
  keepspaces=true
}

% ---------------------------------------------------------
% Ambiente de código quebrável
% Uso: \begin{codigo}{Título} ... \end{codigo}
% ---------------------------------------------------------
\newtcblisting[auto counter, number within=chapter]{codigo}[2][]{
  breakable,
  enhanced,
  listing only,
  listing options={style=codigo},
  colback=codebg,
  colframe=codeframe,
  boxrule=0.4pt,
  arc=2pt,
  left=6mm,right=2mm,top=2mm,bottom=2mm,
  title={Listagem~\thetcbcounter: #2},
  fonttitle=\bfseries,
  #1
}

%=======================================================================
% 19. CAIXA PARA PROMPTS CIENTÍFICAS / ORIENTAÇÕES IA (VERBATIM)
%=======================================================================
\newtcblisting{promptcientifico}[2][]{%
  breakable,
  enhanced,
  listing only,
  colback=gray!10,
  colframe=black,
  boxrule=0.6pt,
  arc=2pt,
  left=4mm,
  right=4mm,
  top=2mm,
  bottom=2mm,
  title={#2},
  fonttitle=\bfseries,
  listing options={
    basicstyle=\ttfamily\small,
    breaklines=true,
    columns=fullflexible,
    keepspaces=true,
    showstringspaces=false
  },
  #1
}



%=======================================================================
% 20. HYPERREF
%=======================================================================
% hyperref: Cria ligações e marcadores no PDF. Conservam-se as opções originais; o
% último colorlinks=true prevalece sobre colorlinks=false.
\usepackage[%
  bookmarks=true,
  citecolor=black,
  filecolor=black,
  linkcolor=black,
  urlcolor=blue,     % Cor específica para URLs
  unicode=false,
  colorlinks=false,
  plainpages=false,
  pdfpagelabels,
  colorlinks=true,  % Ativa cores em vez de caixas
  xetex]{hyperref}

%=======================================================================
% 21. TOPO E RODAPÉ DE PÁGINAS (LOGOS E TEXTO)
%=======================================================================
% eso-pic: Adiciona elementos ao fundo ou ao primeiro plano das páginas no momento de
% as enviar para o PDF.
\usepackage{eso-pic}

\AddToShipoutPictureBG{%           % Início do bloco que adiciona elementos ao fundo da página
  \ifnum\value{page}>2           % Condição: só aplica a partir da página 3
    \AtPageUpperLeft{%            % Posiciona tudo relativamente ao canto superior esquerdo da página
      \begin{tikzpicture}[remember picture, overlay]% Início do ambiente TikZ sobreposto (overlay) e lembrando posição
        % Linhas superior e inferior
        \draw[line width=0.5pt, color=black] (2.5cm, -3cm) -- (\paperwidth - 2.5cm, -3cm); % Linha horizontal superior
        \draw[line width=0.5pt, color=black] (2.5cm, -\paperheight + 3cm) -- (\paperwidth - 2.5cm, -\paperheight + 3cm); % Linha horizontal inferior
        
        % Conteúdo página ímpar
        \ifodd\value{page}           % Verifica se a página é ímpar
          \node [anchor=west, text depth=0pt] at (2.5cm, -2.7cm){\tiny\href{mailto:luisav.ferreira@gmail.com}{Grupo Disciplinar 510}}; % Texto canto superior esquerdo
          \node [anchor=north east] at (\paperwidth-2.5cm, -2.2cm){\includegraphics[width=0.05\paperwidth]{logos/aeB05.pdf}}; % Logo canto superior direito
          \node [anchor=west, text height=0pt] at (2.5cm, -\paperheight + 2.7cm){\tiny\href{mailto:luisav.ferreira@gmail.com}{Luís A.V. Ferreira}}; % Texto canto inferior esquerdo
          \node [anchor=west] at (2.5cm, -\paperheight + 2.5cm){\tiny\jobname.tex \smash{\rule[-0.5ex]{0.6pt}{1ex}}}; % Nome do ficheiro e linha decorativa
        % Conteúdo página par
        \else                        % Caso seja página par
          \node [anchor=west, text depth=0pt] at (2.5cm, -2.7cm){\includegraphics[width=0.05\paperwidth]{logos/aeB05.pdf}}; % Logo canto superior esquerdo
          \node [anchor=east, text depth=0pt] at (\paperwidth - 2.5cm, -2.7cm){\tiny\href{mailto:luisav.ferreira@gmail.com}{Grupo Disciplinar 510}}; % Texto canto superior direito
          \node [anchor=east, text height=0pt] at (\paperwidth - 2.5cm, -\paperheight + 2.7cm){\tiny\href{mailto:luisav.ferreira@gmail.com}{Luís A.V. Ferreira}}; % Texto canto inferior direito
          \node [anchor=east] at (\paperwidth - 2.5cm, -\paperheight + 2.5cm){\tiny\jobname.tex \smash{\rule[-0.5ex]{0.6pt}{1ex}}}; % Nome do ficheiro e linha decorativa
        \fi
      \end{tikzpicture}
    }
  \fi
}



\begin{document}
%%%%%%%%%%%%%%%%%%%%%%%%%%%%%%%%%

\begin{titlepage}
	\centering
	\includegraphics[width=0.15\textwidth]{logos/aeB06.pdf}  \par
	\vspace{1cm}
	{\scshape\LARGE Escola Secundaria José Gomes Ferreira    \par}
	\vspace{1cm}
	{\scshape\Large Física 12º     \par}
	\vspace{1.5cm}
	{\huge\bfseries Modelos de relatório}\par 
	\vspace{2cm}
	{\Large\itshape Nome do aluno}\par
	{\Large\itshape Nº e Turma}\par 
	\vfill
% Bottom of the page
	{\large \today\par}

\begin{tikzpicture}[remember picture, overlay]
\draw[line width = 2pt, color=BLACK ] ($(current page.north west) + (.5in,-.5in)$) rectangle ($(current page.south east) + (-.5in,.5in)$);
\draw[line width = 1pt, color=ELECTRICRED ] ($(current page.north west) + (.55in,-.55in)$) rectangle ($(current page.south east) + (-.55in,.55in)$);
\draw[line width = 1pt, color=TANGERINYELLOW ] ($(current page.north west) + (.6in,-.6in)$) rectangle ($(current page.south east) + (-.6in,.6in)$);
\end{tikzpicture}

\end{titlepage}

\cleardoublepage

%%%%%%%%%%%%%%%%%%%%%%%%%%%%%%%%%%%%%%%%%%%%%%%%

\begin{abstract}
Este espaço destina-se a um resumo breve do trabalho, a redigir no final.
Indicar o objectivo do projecto, a metodologia geral seguida e os principais
resultados obtidos (sem detalhes técnicos).
Manter entre 3 e 6 linhas.
Substituir este texto após concluir integralmente o relatório.
\end{abstract}
\cleardoublepage

%%%%%%%%%%%%%%%%%%%%%%%%%%%%%%%%%%%%%%%%%%%%%%%%

\tableofcontents
\cleardoublepage

%%%%%%%%%%%%%%%%%%%%%%%%%%%%%%%%%%%%%%%%%%%%%%%%

O objectivo deste trabalho consiste em desenvolver um projecto com
componente conceptual e computacional, culminando na produção de
resultados visuais e de ficheiros associados.

A utilização de sistemas de Inteligência Artificial deverá ocorrer
exclusivamente como instrumento de apoio organizacional, linguístico
e estrutural, não podendo substituir o pensamento científico,
a análise crítica nem o trabalho intelectual do aluno.

\vspace{0.5cm}

\textbf{Instruções gerais}

O aluno deverá copiar integralmente o texto apresentado no bloco seguinte
e utilizá-lo como \emph{prompt} num sistema de Inteligência Artificial.

A resposta obtida deverá ser entendida apenas como proposta inicial de
organização do relatório, devendo posteriormente ser analisada,
revista e adaptada criticamente.

\vspace{0.5cm}

%%%%%%%%%%%%%%%%%%%%%%%%%%%%%%%%%%%%%%%%%%%%%%%%

\begin{promptcientifico}{Prompt de Planeamento Científico}

Actua como editor de LaTeX (XeLaTeX) e GERADOR DE ARQUITECTURA DOCUMENTAL
para um trabalho final de Física (12.º), em português europeu pt-PT.

%%%%%%%%%%%%%%%%%%%%%%%%%%%%%%%%%%%%%%%%%%%%%%%%
\textbf{REGRA 0 — MODO ARQUITECTO (ABSOLUTA)}
%%%%%%%%%%%%%%%%%%%%%%%%%%%%%%%%%%%%%%%%%%%%%%%%

A tua saída é apenas um ESQUELETO de relatório.
Não és professor, nem manual, nem explicador.
É proibido produzir texto que ensine, explique, descreva ou contextualize ciência.

Sempre que surja qualquer frase com vocabulário científico/técnico,
substitui-a por placeholders do tipo:
[CONCEITO A EXPLICAR], [EQUAÇÃO A INSERIR PELO ALUNO], [FIGURA A INSERIR], [REFERÊNCIA].

%%%%%%%%%%%%%%%%%%%%%%%%%%%%%%%%%%%%%%%%%%%%%%%%
\textbf{Tarefa}
%%%%%%%%%%%%%%%%%%%%%%%%%%%%%%%%%%%%%%%%%%%%%%%%

Reescreve integralmente um ficheiro \texttt{.tex}, a partir de um MODELO fornecido,
transformando-o num “modelo inteligente” pronto a ser preenchido por um aluno.

%%%%%%%%%%%%%%%%%%%%%%%%%%%%%%%%%%%%%%%%%%%%%%%%
\textbf{Formato da resposta}
%%%%%%%%%%%%%%%%%%%%%%%%%%%%%%%%%%%%%%%%%%%%%%%%

1) Deves devolver APENAS o código LaTeX completo do ficheiro final (um único bloco \texttt{.tex}),
sem qualquer texto fora do código.

2) Mantém exactamente, sem alterações e pela mesma ordem, estas linhas no início do ficheiro:
\begin{itemize}
\item \% !TEX TS-program = xelatex
\item \% pt-PT, AO45, UTF-8
\item \texttt{\textbackslash lccode`\textbackslash-=`\textbackslash-}
\item \texttt{\textbackslash documentclass[12pt,a4paper,twoside]\{report\}}
\item \texttt{\textbackslash input\{auxiliares/preambulo.tex\}}
\end{itemize}

3) Mantém a capa: \texttt{\textbackslash input\{auxiliares/capa.tex\}}.

4) O abstract deve existir e conter apenas 3--6 linhas de texto-modelo genérico,
sem qualquer termo científico.

5) O índice (\texttt{\textbackslash tableofcontents}) deve permanecer logo após o abstract.

%%%%%%%%%%%%%%%%%%%%%%%%%%%%%%%%%%%%%%%%%%%%%%%%
\textbf{Capítulo obrigatório: Orientação Inteligente do Trabalho}
%%%%%%%%%%%%%%%%%%%%%%%%%%%%%%%%%%%%%%%%%%%%%%%%

Deve existir um capítulo com o título exacto:
\begin{center}
\textbf{Orientação Inteligente do Trabalho}
\end{center}

Este capítulo pode conter texto corrido, mas APENAS sobre:
\begin{itemize}
\item ética e honestidade académica;
\item regras de uso de IA (apoio estrutural/linguístico, nunca substituição);
\item organização do trabalho e entregáveis;
\item checklist mínima de conformidade.
\end{itemize}

PROIBIÇÃO: mesmo neste capítulo é proibido usar termos científicos/técnicos do domínio.
Se precisares de os referir, usa apenas [CONCEITO A EXPLICAR] ou [REFERÊNCIA].

%%%%%%%%%%%%%%%%%%%%%%%%%%%%%%%%%%%%%%%%%%%%%%%%
\textbf{Gramática estrutural (OBRIGATÓRIA)}
%%%%%%%%%%%%%%%%%%%%%%%%%%%%%%%%%%%%%%%%%%%%%%%%

1) Usa capítulos numerados (\texttt{\textbackslash chapter}) para a estrutura principal.

2) Dentro de cada \texttt{\textbackslash chapter}, é obrigatório usar \texttt{\textbackslash section}.
É PROIBIDO usar \texttt{\textbackslash subsection} directamente após \texttt{\textbackslash chapter}.

3) É PROIBIDO usar: \texttt{\textbackslash chapter*}, \texttt{\textbackslash section*},
\texttt{\textbackslash subsection*}, \texttt{\textbackslash paragraph}, \texttt{\textbackslash subparagraph}.

%%%%%%%%%%%%%%%%%%%%%%%%%%%%%%%%%%%%%%%%%%%%%%%%
\textbf{REGRA 0.1 — ITENS COM CONTEÚDO CONTROLADO (OBRIGATÓRIA)}
%%%%%%%%%%%%%%%%%%%%%%%%%%%%%%%%%%%%%%%%%%%%%%%%

Em todas as listas \texttt{itemize} (excepto na checklist do capítulo “Orientação Inteligente do Trabalho”):

\begin{itemize}
\item o 1.º \texttt{\textbackslash item} pode ser apenas o rótulo do bloco:
      \texttt{\textbackslash item \textbackslash textbf\{O que escrever aqui\}}
      ou
      \texttt{\textbackslash item \textbackslash textbf\{Perguntas-guia\}};
\item todos os itens seguintes DEVEM ter a forma:
\end{itemize}

\begin{verbatim}
\item <VERBO NO INFINITIVO> [PLACEHOLDER].
\end{verbatim}

É PROIBIDO escrever substantivos técnicos, descrições, ou frases com conteúdo.
Só são permitidas palavras funcionais (artigos, preposições) e placeholders.

Placeholders permitidos:
[CONCEITO A EXPLICAR], [DECISÃO A JUSTIFICAR], [PROCEDIMENTO A DESCREVER],
[PARÂMETRO A DEFINIR], [FICHEIRO A REFERIR], [FERRAMENTA A INDICAR],
[FIGURA A INSERIR], [TABELA A INSERIR], [REFERÊNCIA], [EQUAÇÃO A INSERIR PELO ALUNO].

%%%%%%%%%%%%%%%%%%%%%%%%%%%%%%%%%%%%%%%%%%%%%%%%
\textbf{REGRA 0.2 — LISTA BRANCA DE VERBOS PERMITIDOS}
%%%%%%%%%%%%%%%%%%%%%%%%%%%%%%%%%%%%%%%%%%%%%%%%

Só podes iniciar itens com estes verbos (ou muito próximos):
Definir, Indicar, Inserir, Enumerar, Registar, Referir, Organizar, Justificar, Descrever, Documentar, Verificar.

Se precisares de outra formulação, usa um destes verbos e um placeholder.

%%%%%%%%%%%%%%%%%%%%%%%%%%%%%%%%%%%%%%%%%%%%%%%%
\textbf{REGRA 0.3 — VOCABULÁRIO TÉCNICO PROIBIDO (GERAL)}
%%%%%%%%%%%%%%%%%%%%%%%%%%%%%%%%%%%%%%%%%%%%%%%%

É proibido usar palavras como (ou equivalentes):
algoritmo, algoritmo(s), cálculo, equação, variável, parâmetro(s),
resolução, domínio, discretização, contorno, convergência,
erro(s), propagação, validação, renderização, iluminação, câmara,
malha, vértices, faces, normais, simulação, modelo, teoria, dados.

Se surgirem, substituir a expressão por um placeholder apropriado
(p.ex. [PROCEDIMENTO A DESCREVER], [PARÂMETRO A DEFINIR], [REFERÊNCIA]).

Nota: esta proibição NÃO se aplica a palavras que ocorram dentro de placeholders entre [ ].

%%%%%%%%%%%%%%%%%%%%%%%%%%%%%%%%%%%%%%%%%%%%%%%%
\textbf{Regra sintáctica absoluta por secção (SEM EXCEPÇÕES)}
%%%%%%%%%%%%%%%%%%%%%%%%%%%%%%%%%%%%%%%%%%%%%%%%

Depois de CADA \texttt{\textbackslash section}, \texttt{\textbackslash subsection} e \texttt{\textbackslash subsubsection},
deves escrever EXACTAMENTE dois ambientes \texttt{itemize}, e nada mais:

\begin{verbatim}
\begin{itemize}
\item \textbf{O que escrever aqui}
\item ...
\end{itemize}

\begin{itemize}
\item \textbf{Perguntas-guia}
\item ...
\end{itemize}
\end{verbatim}

É proibido qualquer parágrafo fora destes dois blocos.

%%%%%%%%%%%%%%%%%%%%%%%%%%%%%%%%%%%%%%%%%%%%%%%%
\textbf{Vocabulário proibido (LISTA NEGRA)}
%%%%%%%%%%%%%%%%%%%%%%%%%%%%%%%%%%%%%%%%%%%%%%%%

É proibido escrever qualquer um dos termos (ou equivalentes) abaixo.
Se aparecerem, substitui a frase por placeholders:

\begin{itemize}
\item mecânica quântica; quântico; átomo; hidrogenoide; orbital; orbitais;
\item equação; Schrödinger; função de onda; densidade de probabilidade;
\item simetria; nós; nodos; estados; números quânticos;
\item modelo teórico; teoria; dados científicos; fundamentação teórica;
\item simulação; interpretação física; contexto histórico; importância na física.
\end{itemize}

Nota: esta proibição NÃO se aplica a palavras que ocorram dentro de placeholders entre [ ].

%%%%%%%%%%%%%%%%%%%%%%%%%%%%%%%%%%%%%%%%%%%%%%%%
\textbf{Conteúdo permitido}
%%%%%%%%%%%%%%%%%%%%%%%%%%%%%%%%%%%%%%%%%%%%%%%%

Só é permitido:
\begin{itemize}
\item arquitectura do documento;
\item instruções genéricas de escrita (“Definir…”, “Documentar…”, “Justificar…”),
desde que NÃO incluam vocabulário científico;
\item placeholders.
\end{itemize}

%%%%%%%%%%%%%%%%%%%%%%%%%%%%%%%%%%%%%%%%%%%%%%%%
\textbf{Capítulos obrigatórios (apenas como títulos/arquitectura)}
%%%%%%%%%%%%%%%%%%%%%%%%%%%%%%%%%%%%%%%%%%%%%%%%

Incluir obrigatoriamente capítulos com secções e subsecções adequadas:
\begin{itemize}
\item Metodologia Computacional (grelha/volume; cálculo; iso-valor; [REFERÊNCIA]; exportação OBJ/PLY; Blender; reprodutibilidade);
\item Validação e Análise Crítica (verificações; comparação com [REFERÊNCIA]; limitações; fontes de erro; melhorias);
\item Resultados (figuras PNG; legendas; referência a ficheiros OBJ/PLY);
\item Conclusões;
\item Apêndices: Código (com ambiente \texttt{codigo}), Parâmetros e execuções, Notas adicionais.
\end{itemize}

%%%%%%%%%%%%%%%%%%%%%%%%%%%%%%%%%%%%%%%%%%%%%%%%
\textbf{Código}
%%%%%%%%%%%%%%%%%%%%%%%%%%%%%%%%%%%%%%%%%%%%%%%%

Se inserires código, usa apenas:
\begin{verbatim}
\begin{codigo}{Título}
...
\end{codigo}
\end{verbatim}

%%%%%%%%%%%%%%%%%%%%%%%%%%%%%%%%%%%%%%%%%%%%%%%%
\textbf{Final obrigatório}
%%%%%%%%%%%%%%%%%%%%%%%%%%%%%%%%%%%%%%%%%%%%%%%%

No fim do ficheiro, manter:
\begin{verbatim}
\printbibliography[title={Referências}]
\label{fim}
\end{document}
\end{verbatim}

E acrescentar:
\begin{verbatim}
% CHECK: sem equações, sem símbolos gregos,
% sem parágrafos fora de listas (excepto no cap. Orientação),
% sem conteúdo científico, sem termos técnicos do domínio.
\end{verbatim}

%%%%%%%%%%%%%%%%%%%%%%%%%%%%%%%%%%%%%%%%%%%%%%%%

Agora segue o MODELO ORIGINAL:

[COLAR AQUI O FICHEIRO .TEX ORIGINAL INTEIRO]

Reescreve integralmente e devolve apenas o ficheiro final.

\end{promptcientifico}

\vspace{0.5cm}

\textbf{Nota importante}

A estrutura gerada pela Inteligência Artificial não constitui
uma versão definitiva do trabalho.

O aluno deverá analisá-la criticamente, adaptá-la, justificar as alterações
efectuadas e assumir plena responsabilidade científica pelo resultado final.

O objectivo fundamental consiste em aprender a organizar pensamento científico
autónomo, e não delegar a sua construção.

\vspace{0.5cm}

\textbf{Orientações para investigação e construção científica do trabalho}

Após a obtenção da estrutura inicial do relatório, o aluno deverá
prosseguir o desenvolvimento do trabalho através de investigação
bibliográfica autónoma.

A Inteligência Artificial fornece apenas uma organização preliminar
do documento. O conhecimento científico necessário à sua elaboração
deve ser obtido a partir de fontes académicas e bibliográficas
credíveis.

Recomenda-se a consulta regular de repositórios digitais de acesso
aberto, tais como:

\begin{itemize}
\item \url{https://archive.org/}
\item \url{https://arxiv.org/}
\item \url{https://scholar.google.com/}
\item \url{https://www.doabooks.org/}
\item \url{https://www.gutenberg.org/}
\item \url{https://sci-hub.red/}
\item bibliotecas universitárias e repositórios institucionais.
\end{itemize}

Sempre que possível, o aluno deverá obter os documentos em formato
PDF para leitura integral e, quando disponível, versões em formato
texto, permitindo pesquisa directa de conceitos, definições e
terminologia relevante.

A bibliografia deve ser construída progressivamente ao longo do
desenvolvimento do trabalho, e não apenas na fase final.
Cada nova fonte consultada deve ser imediatamente registada no
ficheiro bibliográfico correspondente.

As referências bibliográficas devem ser citadas no momento em que
são utilizadas no texto, justificando afirmações, definições ou
interpretações apresentadas. A inclusão tardia de referências apenas
no final do relatório constitui uma prática incorrecta e deverá ser
evitada.

O relatório científico deve reflectir explicitamente o percurso de
investigação realizado, evidenciando a relação entre leitura,
compreensão, escrita e fundamentação bibliográfica.


%%%%%%%%%%%%%%%%%%%%%%%%%%%%%%%%%%%%%%%%%%%%%%%%

\printbibliography[title={Referências}]

\label{fim}
\end{document}
